\section{Introduction}
\label{sec:intro}

Council tax is the United Kingdom's principal recurrent tax on residential property, with gross receipts of approximately \pounds 62 billion forecast for 2026--27 \citep{obr2025}. Its design has attracted sustained criticism in the public-finance literature \citep{mirrlees2011, adam2020}. Dwellings in England and Scotland remain assessed on their estimated value as at 1 April 1991; the banding structure is capped, with the top band in England (Band H) covering all properties worth \pounds 320{,}000 or more \emph{in 1991 prices}, so that properties of widely differing current values share the top band and face nearly identical bills; and effective tax rates decline as a share of property value as that value rises. Council tax is therefore regressive with respect to the value of the asset it nominally taxes, bears a weak relationship to households' current ability to pay, and discourages investment in property improvements.

This paper examines the household-level consequences of abolishing council tax and replacing its net revenue, pound for pound, with a flat annual tax on the unimproved value of land. The land value tax (LVT) has a long intellectual pedigree, from \citet{george1879} through the optimal-taxation literature \citep{ramsey1927} to the review by \citet{mirrlees2011}, resting on the observation that the supply of land is fixed: a tax on pure land rents raises revenue without discouraging productive activity. The UK policy debate, however, has to date lacked a household-level quantification of the gains and losses such a reform would generate.

We provide such a quantification. Using PolicyEngine UK, an open-source static microsimulation model of the UK tax and benefit system built on an enhanced version of the Family Resources Survey \citep{policyengine}, we proceed in three steps. First, we impute a land value to every household by combining survey property wealth with region-specific land shares and an allocation of corporate-sector land, calibrating national aggregates to the ONS National Balance Sheet \citep{ons2025}. Second, we reconstruct council tax liabilities bottom-up by region, band and discount status, calibrated to Valuation Office Agency dwelling counts. Third, we simulate replacing the model's net council tax revenue of \pounds 57.6 billion (gross liabilities of \pounds 61.9 billion less \pounds 4.3 billion of council tax reduction) with a flat-rate LVT for fiscal year 2026--27; Section~\ref{sec:methodology} discusses how this model-based figure compares with the OBR receipts forecast. Against an estimated UK land base of \pounds 7.46 trillion, we calculate that revenue neutrality requires a rate of 0.77 per cent of land value per year.

Our results suggest three principal findings. First, the swap is progressive across most of the income distribution: we estimate that the average household in each of income deciles 1--8 gains, with the poorest decile gaining \pounds 481 per year --- 2.3 per cent of its net income --- while the top two deciles lose \pounds 991 and \pounds 966 per year respectively; 68 per cent of households gain overall. Second, the reform reduces poverty: the absolute poverty rate before housing costs falls by 0.65 percentage points, and by 1.04 points after housing costs. Third, the Gini coefficient of equivalised household net income is essentially unchanged (+0.0004) even though the burden shifts markedly towards the wealthy: ranked by household wealth rather than income, the top wealth decile loses \pounds 5,573 per year on average and decile 9 loses \pounds 1,042, while wealth deciles 1--7 gain between \pounds 484 and \pounds 1,464. These findings can be reconciled by observing that land wealth is far more unequally distributed than income --- we estimate that the top wealth decile holds 48.5 per cent of household land --- but is only weakly correlated with current income, so that a tax on land redistributes across the wealth distribution in ways that income-based inequality statistics register only faintly.

In addition to the central budget-neutral scenario, we trace out a rate grid from 0.5 to 5 per cent, documenting how revenue (from \pounds 37.3 billion to \pounds 373.2 billion), poverty and inequality vary with the rate, and we decompose the tax base into household and corporate land.

The remainder of the paper proceeds as follows. Section~\ref{sec:literature} situates the analysis in the literature on land taxation and UK property-tax reform. Section~\ref{sec:methodology} describes the data, the land-value imputation, the council tax baseline and the reform implementation. Section~\ref{sec:results} presents the distributional results. Section~\ref{sec:discussion} summarises the findings, discusses limitations and policy implications, and concludes. An appendix reports the full rate grid, a summary of landless households, a worked single-household example and a data and parameter summary.
